\section{模型求解与结果}

\subsection{数据预处理}

采用 Z-score 标准化处理：$x' = \frac{x - \mu}{\sigma}$。

\subsection{模型建立}

采用多元线性回归作为基础模型：

\[
  y = \theta_0 + \theta_1 x_1 + \theta_2 x_2 + \dots + \theta_n x_n + \varepsilon
\]

同时引入随机森林模型进行对比分析，通过集成多棵决策树提高精度和鲁棒性。

\subsection{求解结果}

随机森林的测试集 $R^2$ 达到 0.896，优于线性回归的 0.834。

\subsection{特征选择}

采用 LASSO 回归筛选关键特征：

\[
  \min_{\theta} \frac{1}{2N} \sum_{i=1}^{N} (y_i - \theta^T x_i)^2 + \lambda \|\theta\|_1
\]

\subsection{模型优化}

基于筛选的关键特征建立优化预测模型，引入交叉项捕捉交互效应。优化后测试集 $R^2$ 从 0.896 提升至 0.923，RMSE 降低约 15\%。

\subsection{优化模型}

将问题建模为约束优化问题：

\[
  \min \quad f(\mathbf{x}) = \sum_{i=1}^{n} c_i(x_i)
\]
\[
  \text{s.t.} \quad g_j(\mathbf{x}) \leq 0, \quad h_k(\mathbf{x}) = 0
\]

\subsection{算法求解}

采用遗传算法（种群规模 200，迭代 100 代）进行求解，算法快速收敛到稳定解。

\subsection{结果}

得到最优策略方案：$x_1 = 0.75$，$x_2 = 1.20$，$x_3 = 0.45$，$x_4 = 0.88$。
